\documentclass[11pt]{article}
\usepackage[utf8]{inputenc}
\usepackage[T1]{fontenc}
\usepackage[margin=1.8cm]{geometry}
\usepackage{amsmath,amssymb}
\usepackage{array}
\setlength{\parindent}{0pt}
\setlength{\parskip}{5pt}
\newcommand{\work}[1]{\vspace{#1}}
\begin{document}

\begin{center}
{\Large\bfseries GEOMETRÍA 8 — WORKSHOP DE CONTINUACIÓN}\\[3pt]
{\large Áreas sombreadas, escala, área y perímetro integrado}\\[3pt]
Semanas 4–7 del tercer periodo 2026
\end{center}

\hrule
\vspace{6pt}
\textbf{Método obligatorio:} dibuja la figura, identifica las piezas, escribe la expresión completa, calcula y verifica las unidades.
\vspace{6pt}
\hrule

\section*{A. Semana 4 — Áreas sombreadas simples}

\textbf{1.} Un cuadrado de lado $10\,\mathrm{cm}$ contiene un círculo interior de radio $3\,\mathrm{cm}$. Calcula el área del cuadrado que queda fuera del círculo.
\work{2.0cm}

\textbf{2.} Un rectángulo de $15\,\mathrm{cm}\times 8\,\mathrm{cm}$ contiene un semicírculo de radio $4\,\mathrm{cm}$. Calcula el área que queda al retirar el semicírculo.
\work{2.0cm}

\textbf{3.} Un cuadrado de lado $12\,\mathrm{cm}$ contiene un cuadrante de radio $12\,\mathrm{cm}$. Calcula el área que queda fuera del cuadrante.
\work{2.0cm}

\textbf{4.} Un círculo de radio $7\,\mathrm{cm}$ contiene una región cuadrada de lado $8\,\mathrm{cm}$. Calcula la diferencia entre el área del círculo y el área del cuadrado.
\work{2.0cm}

\newpage
\section*{B. Semana 5 — Áreas sombreadas complejas}

\textbf{5.} Una corona circular tiene radio exterior $10\,\mathrm{cm}$ y radio interior $6\,\mathrm{cm}$. Calcula su área.
\work{2.1cm}

\textbf{6.} Un cuadrado de lado $16\,\mathrm{cm}$ tiene cuatro cuadrantes de radio $8\,\mathrm{cm}$, uno en cada esquina. Los cuatro cuadrantes equivalen a un círculo completo. Calcula el área central restante.
\work{2.1cm}

\textbf{7.} Un rectángulo de $18\,\mathrm{cm}\times10\,\mathrm{cm}$ tiene dos huecos circulares iguales de radio $2\,\mathrm{cm}$. Calcula el área restante.
\work{2.1cm}

\textbf{8.} Un círculo de radio $8\,\mathrm{cm}$ tiene dos huecos circulares de radio $2\,\mathrm{cm}$. Calcula el área que queda.
\work{2.1cm}

\textbf{9.} Un rectángulo de $24\,\mathrm{cm}\times12\,\mathrm{cm}$ pierde dos semicírculos de radio $6\,\mathrm{cm}$. Calcula el área restante.
\work{2.1cm}

\newpage
\section*{C. Semana 6 — Perímetro vs. área bajo escala}

\textbf{10.} Un cuadrado cambia de lado $5\,\mathrm{cm}$ a lado $15\,\mathrm{cm}$. Determina el factor de escala $k$ y explica por qué factor cambian el perímetro y el área.
\work{2.3cm}

\textbf{11.} Un círculo cambia de radio $4\,\mathrm{cm}$ a radio $2\,\mathrm{cm}$. Determina $k$ y explica por qué factor cambian la circunferencia y el área.
\work{2.3cm}

\textbf{12.} Una figura semejante tiene factor de escala $k=1.5$ y área original $40\,\mathrm{cm}^2$. Calcula el área nueva.
\work{2.3cm}

\textbf{13.} Una figura tiene perímetro original $36\,\mathrm{cm}$ y área original $50\,\mathrm{cm}^2$. Se amplía con $k=2.5$. Calcula el nuevo perímetro y la nueva área.
\work{2.3cm}

\newpage
\section*{D. Semana 7 — Reto integrado de área y perímetro}

\textbf{14.} Una pista tipo estadio tiene dos tramos rectos de $24\,\mathrm{m}$ y dos extremos semicirculares de radio $5\,\mathrm{m}$. Calcula el área total y el perímetro exterior.
\work{2.5cm}

\textbf{15.} Un patio rectangular de $18\,\mathrm{m}\times14\,\mathrm{m}$ contiene una fuente circular de radio $4\,\mathrm{m}$. Calcula la superficie libre del patio y el perímetro exterior del rectángulo.
\work{2.5cm}

\textbf{16.} La figura del ejercicio 15 se amplía de manera semejante con $k=2.5$. Sin reconstruir todas las fórmulas, explica por qué factor cambia el perímetro y por qué factor cambia el área.
\work{2.5cm}

\textbf{17.} Un parque circular tiene radio exterior $9\,\mathrm{m}$ y un jardín circular interior de radio $4\,\mathrm{m}$. Calcula el área de la corona y la longitud total de ambos bordes circulares.
\work{2.5cm}

\section*{Exit ticket}
\textbf{a)} ¿Qué palabras en un problema te indican que debes calcular perímetro?\\
\textbf{b)} ¿Qué palabras te indican que debes calcular área?\\
\textbf{c)} Si una figura se amplía con factor $k=4$, ¿por qué factor cambia su área?

\end{document}
